%%==================================================================
%% IEEE Transactions on Medical Imaging (TMI) -- LaTeX Template
%%==================================================================
%% This template is adapted from the official IEEEtran bare_jrnl.tex
%% (IEEEtran.cls v1.8b, downloaded from CTAN:
%%  https://ctan.org/pkg/ieeetran). TMI uses the standard IEEEtran
%% class in journal mode; there is no TMI-specific .cls file.
%%
%% TMI-specific notes (per https://ieeetmi.org/authors-instructions/,
%% Revision 10.2, Oct 9 2025):
%%   * Initial submission MUST be <= 10 pages including references.
%%   * Double-column, single-spaced, justified, 10pt font.
%%   * Abstract < 250 words, plus index terms.
%%   * Do NOT include author biographies in initial submission.
%%   * Figures/tables must be inline in the main text (not at end).
%%   * References use IEEEtran.bst format with author-name style.
%%   * Cover letter strongly discouraged unless extending a conference
%%     paper or disclosing prior review.
%%   * Graphical abstract optional, uploaded as supporting document.
%%   * Open-source code strongly encouraged (link in submission).
%%
%% Compile with:  pdflatex TMI_template && bibtex TMI_template &&
%%                pdflatex TMI_template && pdflatex TMI_template
%%==================================================================

\documentclass[journal,10pt]{IEEEtran}

% ---- Common packages for medical imaging papers ----
% Adjust to your local TeX Live installation; algorithm/algorithmic
% may need `texlive-science` (Debian/Ubuntu) or `algorithm2e`/`algpseudocode`
% as alternatives. They are commented out here to keep the skeleton
% compilable on a minimal TeX Live install.
\usepackage{amsmath,amssymb,amsfonts}
\usepackage{graphicx}
\usepackage{booktabs}
\usepackage{multirow}
\usepackage{cite}
\usepackage{url}
\usepackage[hidelinks]{hyperref}
% \usepackage{algorithm}
% \usepackage{algorithmic}

% ---- TMI-specific hygiene ----
% TMI does not publish author biographies in the initial submission.
% Do not invoke \IEEEbiography until requested at acceptance stage.

\begin{document}

\title{Rectified Flow for Chromosome Detection: Stable Coupling and Few-Step Inference}

\author{
  \IEEEauthorblockN{First~Author,~Second~Author,~and~Senior~Author~\IEEEmembership{Member,~IEEE}}
  \IEEEauthorblockA{Department of X, University of Y\\
  City, Country\\
  Email: \{first, second, senior\}@university.edu}
}

\markboth{IEEE Transactions on Medical Imaging,~Vol.~XX, No.~X, Month~2026}{Author \MakeLowercase{\textit{et al.}}: Rectified Flow for Chromosome Detection}

\maketitle

\begin{abstract}
% TMI requires abstract < 250 words. Frame: clinical problem (karyotyping
% bottleneck), methodological gap (diffusion-based detectors suffer
% training instability), proposed contribution (Rectified Flow with
% Stochastic Coupling + DPM-Solver++ 2-step inference + Top-K pruning),
% and quantitative result (+0.082 mAP on two public chromosome datasets).
% State clinical implications.
Chromosome karyotyping is the first-line diagnostic test for hematological
malignancies and congenital syndromes, yet manual analysis requires hours
per specimen. Diffusion-based object detectors offer a probabilistic
framework for chromosome detection but suffer training instability caused
by optimal-transport (OT) coupling diversity collapse. We propose a
Rectified Flow (RF) detector with Stochastic Coupling, theoretically
grounded by a diversity-collapse analysis, and accelerated by a
DPM-Solver++ two-step sampler with Top-K box pruning. On two public
chromosome datasets, our method improves mAP by 0.082 over
DiffusionDet while reducing inference to 2 steps at $X$ FPS. The
stabilizer eliminates seed-level variance, enabling clinically
deployable, real-time metaphase analysis. Code and trained models
are released to support reproducibility.
\end{abstract}

\begin{IEEEkeywords}
Chromosome detection, rectified flow, diffusion models, optimal transport, few-step inference, karyotyping, medical image analysis.
\end{IEEEkeywords}

% ====================================================================
\section{Introduction}
\label{sec:intro}
% Per TMI editorial (Shan, Kruger, Wang, TMI 44(12), 2025), frame the
% problem around (i) unmet clinical need (karyotyping turnaround time),
% (ii) methodological gap (DiffusionDet training instability), and
% (iii) the paradigm-shifting angle of RF for detection.

\IEEEPARstart{C}{hromosome} analysis via karyotyping is the first-line
test for hematological malignancies and congenital syndromes
\cite{example_ref_2024}. Manual karyotyping takes a skilled
cytogeneticist several hours per specimen, creating a throughput
bottleneck in clinical genetics laboratories. Recent object-detection
frameworks based on denoising diffusion models (e.g.,
DiffusionDet~\cite{diffusiondet}) offer a unified probabilistic
formulation for detection, but their training is sensitive to
coupling diversity and convergence is brittle on dense, small-object
modalities such as metaphase chromosomes.

% ... rest of introduction ...

% ====================================================================
\section{Related Work}
\label{sec:related}

% Briefly cover: (a) classical chromosome detection, (b) diffusion-based
% detectors, (c) Rectified Flow theory, (d) OT coupling literature.
% TMI prefers concise related work; compress aggressively to stay within
% 10-page initial-submission limit (references count toward the cap).

% ====================================================================
\section{Methodology}
\label{sec:method}

\subsection{Rectified Flow for Object Detection}
% Provide formal RF formulation; reference Liu et al. ICLR 2023.

\subsection{Stochastic Coupling: Resolving OT Diversity Collapse}
% State the OT diversity collapse theorem, then introduce the
% Stochastic Coupling stabilizer. TMI's SIER criteria ask for
% mathematical depth: include proposition statements and proof sketches
% in the main text or appendix.

\subsection{DPM-Solver++ 2-Step Inference with Top-K Pruning}
% Describe the few-step sampler and the Top-K box pruning that
% reduces inference cost while preserving AP.

% ====================================================================
\section{Experiments}
\label{sec:exp}

\subsection{Datasets}
% Two public datasets: Chromosome20240904 and 24 Chromosomes Object.
% Report patient-level splits, ethics statement, and IRB status if any.
% TMI requires disclosure of dataset availability.

\subsection{Evaluation Metrics}
% mAP, per-class AP, FPS, statistical significance across seeds.

\subsection{Comparison with State-of-the-Art}
% Compare vs. Cascade R-CNN, YOLOX-S, DiffusionDet.
% TMI reviewers expect at least 3 strong baselines, statistically
% meaningful comparisons (multi-seed), and a discussion of failure modes.

\subsection{Ablation Studies}
% (1) RF vs. standard DDPM coupling, (2) Stochastic Coupling on/off,
% (3) number of inference steps, (4) Top-K threshold.

% ====================================================================
\section{Discussion}
\label{sec:discussion}
% Address clinical translatability, limitations (e.g., overlapping
% chromosomes, low-bandwidth clinical deployment), and broader
% applicability of RF to other detection tasks in medical imaging.

% ====================================================================
\section{Conclusion}
\label{sec:conclusion}

% ====================================================================
\section*{Acknowledgments}
% Funders, data providers, code availability statement.
% TMI strongly encourages open-source code; include GitHub URL.

% ====================================================================
% Bibliography options:
%  (a) If you have a references.bib file, use:
%        \bibliographystyle{IEEEtran}
%        \bibliography{IEEEabrv,references}
%  (b) For initial submission, inline thebibliography is acceptable:
\begin{thebibliography}{99}

\bibitem{example_ref_2024}
J.~Doe and A.~Smith, ``Example reference for chromosome karyotyping,''
\emph{IEEE Trans. Med. Imag.}, vol.~43, no.~5, pp.~1234--1245, May 2024.

\bibitem{diffusiondet}
S.~Jia \emph{et al.}, ``DiffusionDet: Diffusion model for object detection,''
in \emph{Proc. Int. Conf. Learn. Represent. (ICLR)}, 2023.

\end{thebibliography}

\end{document}
